\begin{spacing}{1.15}
\footnotesize
\noindent\textbf{EN01} \quad EN \quad Speed limits \quad easy \\
\textit{Question:} What is the speed limit in a built-up area in Malta? \\
\textit{Reference answer:} In a built-up area (town or village) the default speed limit is 50 km/h unless a road sign indicates otherwise. \\
\textit{Key facts (accepted aliases):} 50 km/h (50km/h, 50 kilometres per hour, 50 kilometers per hour, 50 kph); built-up (built up, urban area, urban areas, towns and villages, town or village, towns or villages, residential) \\
\textit{Expected citation:} S.L. 65.11
\par\vspace{5pt}
\noindent\textbf{EN02} \quad EN \quad Speed limits \quad medium \\
\textit{Question:} How fast can I legally drive on the open road in Malta, and is there a motorway limit? \\
\textit{Reference answer:} Outside towns and villages the open-road limit is 80 km/h. Malta has no motorways, so 80 km/h is the maximum general speed limit; the Authority may set lower signed limits (for example 30 km/h near schools or in residential zones). \\
\textit{Key facts (accepted aliases):} 80 km/h (80km/h, 80 kilometres per hour, 80 kilometers per hour, 80 kph); no motorway (no motorways, not have motorways, not have any motorway, does not have motorways, doesn't have motorways, there are no motorways, has no motorway, without motorways, no highway, no highways, not have highways, no motorway limit, no dedicated motorway) \\
\textit{Expected citation:} S.L. 65.11
\par\vspace{5pt}
\noindent\textbf{EN03} \quad EN \quad Fine payment \quad easy \\
\textit{Question:} What happens to a traffic fine if I do not pay it within 15 days? \\
\textit{Reference answer:} If the fine is not paid within 15 days of issue it doubles, and if still unpaid after 30 days it triples. \\
\textit{Key facts (accepted aliases):} double (doubles, doubled, twice, two times); 15 days (fifteen days, 15-day, 15 day) \\
\textit{Expected citation:} Cap. 65
\par\vspace{5pt}
\noindent\textbf{EN04} \quad EN \quad Fine payment \quad medium \\
\textit{Question:} Where can I pay a contravention online and what happens if I leave it unpaid for a month? \\
\textit{Reference answer:} You can pay online at contraventions.gov.mt within 15 days; the fine doubles after 15 days and triples after 30 days. \\
\textit{Key facts (accepted aliases):} contraventions.gov.mt (); triple (triples, tripled, three times, threefold); 30 days (thirty days, 30-day, 30 day) \\
\textit{Expected citation:} Cap. 65 or Driver FAQ
\par\vspace{5pt}
\noindent\textbf{EN05} \quad EN \quad Alcohol \quad easy \\
\textit{Question:} What is the general blood alcohol limit for driving in Malta? \\
\textit{Reference answer:} The general limit is 50 mg of alcohol per 100 ml of blood (0.05 g/100ml), equivalent to 22 microgrammes per 100 ml of breath. \\
\textit{Key facts (accepted aliases):} 50 mg (50mg, 50 milligram, 50 milligrammes, 50 milligrams); 0.05 (0,05, 0.5 g/l, 0.5g/l, 0.5 grams per litre, 0.5 grammes per litre) \\
\textit{Expected citation:} Cap. 65
\par\vspace{5pt}
\noindent\textbf{EN06} \quad EN \quad Alcohol \quad medium \\
\textit{Question:} Can I legally refuse a breathalyser test when a police officer asks me to take one? \\
\textit{Reference answer:} No. Refusing a breath test without reasonable excuse is an offence treated as equivalent to exceeding the drink-drive limit. \\
\textit{Key facts (accepted aliases):} offence (offense, an offence, criminal, illegal, punishable, against the law, liable to); equivalent (same as, treated the same, treated as, treated in the same way, as if, tantamount, same penalty, same penalties, as though, same consequences) \\
\textit{Expected citation:} Cap. 65
\par\vspace{5pt}
\noindent\textbf{EN07} \quad EN \quad Alcohol \quad hard \\
\textit{Question:} I only got my licence last year. What blood alcohol limit applies to me and why? \\
\textit{Reference answer:} As a new driver on a probationary licence you are subject to the lower limit of 20 mg per 100 ml of blood (0.02 g/100ml), the same reduced limit that applies to commercial drivers. \\
\textit{Key facts (accepted aliases):} 20 mg (20mg, 20 milligram, 20 milligrammes, 20 milligrams); 0.02 (0,02, 0.2 g/l, 0.2g/l, 0.2 grams per litre, 0.2 grammes per litre); probationary (probation, new driver, newly licensed, newly licenced, novice) \\
\textit{Expected citation:} Cap. 65
\par\vspace{5pt}
\noindent\textbf{EN08} \quad EN \quad Penalty points \quad easy \\
\textit{Question:} How many penalty points lead to a licence suspension in Malta, and over what period? \\
\textit{Reference answer:} Accumulating 12 penalty points within any 12-month period leads to revocation of the licence; after a first revocation a new licence can be obtained after two months. \\
\textit{Key facts (accepted aliases):} 12 (twelve); 12-month (12 month, twelve month, twelve-month, 12 months, a year, one year, one-year, any 12, any twelve) \\
\textit{Expected citation:} S.L. 65.18
\par\vspace{5pt}
\noindent\textbf{EN09} \quad EN \quad Penalty points \quad hard \\
\textit{Question:} I am a new driver and I have 7 penalty points. Am I about to lose my licence? \\
\textit{Reference answer:} Not on points alone. New drivers on a probationary licence face revocation at 12 points, the same threshold as full licences, so 7 points is still below the limit (although the lower alcohol limit and other probationary rules still apply). \\
\textit{Key facts (accepted aliases):} 12 (twelve); probationary (probation, new driver, new drivers, newly licensed, novice) \\
\textit{Expected citation:} S.L. 65.18
\par\vspace{5pt}
\noindent\textbf{EN10} \quad EN \quad Accidents \quad medium \\
\textit{Question:} After a minor crash the other driver wants to settle privately. Am I obliged to, or can I involve my insurance? \\
\textit{Reference answer:} You must stop and exchange details, but you are not obliged to settle privately and may involve your insurance. \\
\textit{Key facts (accepted aliases):} stop (stopping, stopped, remain at the scene, stay at the scene); exchange (exchanging, provide your details, give your details, share details, share your details, swap details, provide details); not obliged (not obligated, no obligation, not required, not legally required, under no obligation, do not have to, don't have to, not compelled, not forced, not bound, not mandatory, are not required); insurance (insurer, insurers) \\
\textit{Expected citation:} Cap. 65
\par\vspace{5pt}
\noindent\textbf{EN11} \quad EN \quad Licensing \quad medium \\
\textit{Question:} What is the minimum age to hold a full category B car licence in Malta? \\
\textit{Reference answer:} The minimum age for a full category B licence is 18 (a learner may start at 17). \\
\textit{Key facts (accepted aliases):} 18 (eighteen) \\
\textit{Expected citation:} S.L. 65.18
\par\vspace{5pt}
\noindent\textbf{EN12} \quad EN \quad Micromobility \quad hard \\
\textit{Question:} What are the speed limits for an e-kickscooter on the road compared with on a footpath? \\
\textit{Reference answer:} An e-kickscooter may do up to 20 km/h on roads and 10 km/h on footpaths, promenades and shared paths. \\
\textit{Key facts (accepted aliases):} 20 km/h (20km/h, 20 kilometres per hour, 20 kilometers per hour, 20 kph); 10 km/h (10km/h, 10 kilometres per hour, 10 kilometers per hour, 10 kph); footpath (footpaths, pavement, pavements, sidewalk, sidewalks, promenade, promenades) \\
\textit{Expected citation:} S.L. 65.32
\par\vspace{5pt}
\noindent\textbf{MT01} \quad MT \quad Speed limits \quad easy \\
\textit{Question:} X'inhu l-limitu tal-veloċità f'żona urbana f'Malta? \\
\textit{Reference answer:} F'żona urbana (belt jew raħal) l-limitu hu 50 km/s sakemm ma jkunx indikat mod ieħor b'sinjal. \\
\textit{Key facts (accepted aliases):} 50 (ħamsin) \\
\textit{Expected citation:} S.L. 65.11
\par\vspace{5pt}
\noindent\textbf{MT02} \quad MT \quad Speed limits \quad medium \\
\textit{Question:} Kemm nista' nsuq malajr fuq triq miftuħa, u hemm limitu għall-awtostrada? \\
\textit{Reference answer:} Barra l-bliet u l-irħula l-limitu hu 80 km/s. Malta m'għandhiex awtostradi, mela 80 km/s hu l-limitu ġenerali massimu; l-Awtorità tista' tistabbilixxi limiti aktar baxxi indikati b'sinjali (eż. 30 km/s ħdejn skejjel). \\
\textit{Key facts (accepted aliases):} 80 (tmenin); awtostrad (awtostradi, motorway, motorways, highway) \\
\textit{Expected citation:} S.L. 65.11
\par\vspace{5pt}
\noindent\textbf{MT03} \quad MT \quad Fine payment \quad easy \\
\textit{Question:} X'jiġri mal-multa jekk ma nħallashiex fi żmien 15-il jum? \\
\textit{Reference answer:} Jekk il-multa ma titħallasx fi 15-il jum tirdoppja, u wara 30 jum tittripla. \\
\textit{Key facts (accepted aliases):} doppja (irdoppja, tirdoppja, jirdoppja, rduppjata, irduppjata, darbtejn, doppju, id-doppju, doubles); 15 (ħmistax) \\
\textit{Expected citation:} Cap. 65 or Kap. 65
\par\vspace{5pt}
\noindent\textbf{MT04} \quad MT \quad Fine payment \quad medium \\
\textit{Question:} Fejn nista' nħallas kontravvenzjoni online, u x'jiġri wara 30 jum? \\
\textit{Reference answer:} Tista' tħallas online fuq contraventions.gov.mt; il-multa tirdoppja wara 15-il jum u tittripla wara 30 jum. \\
\textit{Key facts (accepted aliases):} contraventions.gov.mt (); tripla (tittripla, jittripla, tliet darbiet, triplu, it-triplu, ittriplata, triples); 30 (tletin) \\
\textit{Expected citation:} Cap. 65 or Kap. 65 or Driver FAQ
\par\vspace{5pt}
\noindent\textbf{MT05} \quad MT \quad Alcohol \quad easy \\
\textit{Question:} X'inhu l-limitu ġenerali tal-alkoħol fid-demm għas-sewqan? \\
\textit{Reference answer:} Il-limitu ġenerali hu 50 mg ta' alkoħol kull 100 ml demm (0.05 g/100ml), jew 22 mikrogramma kull 100 ml nifs. \\
\textit{Key facts (accepted aliases):} 50 (ħamsin); 0.05 (0,05, 0.5 g/l, 0.5g/l) \\
\textit{Expected citation:} Cap. 65 or Kap. 65
\par\vspace{5pt}
\noindent\textbf{MT06} \quad MT \quad Alcohol \quad medium \\
\textit{Question:} Nista' nirrifjuta test tan-nifs jekk jitolbni l-pulizija? \\
\textit{Reference answer:} Le. Ir-rifjut ta' test tan-nifs bla skuża raġonevoli huwa reat ittrattat bħal qbiż tal-limitu tal-alkoħol. \\
\textit{Key facts (accepted aliases):} reat (offiża, illegali, kontra l-liġi, offence, punibbli, offense, reat kriminali); rifjut (tirrifjuta, jirrifjuta, irrifjuta, rifjuta, tirrifjutax, refus, rrifjuta) \\
\textit{Expected citation:} Cap. 65 or Kap. 65
\par\vspace{5pt}
\noindent\textbf{MT07} \quad MT \quad Alcohol \quad hard \\
\textit{Question:} X'inhu l-limitu tal-alkoħol għal sewwieq professjonali jew ġdid? \\
\textit{Reference answer:} Għal sewwieqa ta' vetturi kummerċjali u sewwieqa ġodda (probatorji) l-limitu hu aktar baxx: 20 mg kull 100 ml (0.02 g/100ml). \\
\textit{Key facts (accepted aliases):} 20 (għoxrin); 0.02 (0,02, 0.2 g/l, 0.2g/l) \\
\textit{Expected citation:} Cap. 65 or Kap. 65
\par\vspace{5pt}
\noindent\textbf{MT08} \quad MT \quad Penalty points \quad easy \\
\textit{Question:} Kemm-il punt tal-penali jwassal għal sospensjoni tal-liċenzja, u f'liema perjodu? \\
\textit{Reference answer:} Tnax-il (12) punt fi kwalunkwe perjodu ta' 12-il xahar iwasslu għal revoka tal-liċenzja; wara l-ewwel revoka tista' tinkiseb liċenzja ġdida wara xahrejn. \\
\textit{Key facts (accepted aliases):} 12 (tnax) \\
\textit{Expected citation:} S.L. 65.18
\par\vspace{5pt}
\noindent\textbf{MT09} \quad MT \quad Penalty points \quad hard \\
\textit{Question:} Jien sewwieq ġdid b'7 punti tal-penali. Qiegħed f'riskju li nitlef il-liċenzja? \\
\textit{Reference answer:} Mhux fuq il-punti biss. Sewwieqa ġodda fuq liċenzja probatorja jiġu revokati b'12-il punt, l-istess limitu bħal liċenzja sħiħa, mela 7 punti għadhom taħt il-limitu (għalkemm il-limitu aktar baxx tal-alkoħol u regoli oħra xorta japplikaw). \\
\textit{Key facts (accepted aliases):} 12 (tnax); probatorja (probatorju, probazzjoni, sewwieq ġdid, sewwieqa ġodda, ġdid, ġodda) \\
\textit{Expected citation:} S.L. 65.18
\par\vspace{5pt}
\noindent\textbf{MT10} \quad MT \quad Accidents \quad medium \\
\textit{Question:} Wara inċident żgħir, obbligat insolvi privatament jew nista' ninvolvi l-assigurazzjoni? \\
\textit{Reference answer:} Trid tieqaf u tiskambja d-dettalji, imma m'intix obbligat issolvi privatament u tista' tinvolvi l-assigurazzjoni. \\
\textit{Key facts (accepted aliases):} tieqaf (tieqafx, waqfa, tibqa' fuq il-post, tieqaf fuq il-post, stop, trid tieqaf); tiskambja (skambju, tagħti d-dettalji, taqsam id-dettalji, tiskambjaw, dettalji, tibdel id-dettalji); assigurazzjoni (assikurazzjoni, insurance, assiguratur) \\
\textit{Expected citation:} Cap. 65 or Kap. 65
\par\vspace{5pt}
\noindent\textbf{MT11} \quad MT \quad Licensing \quad medium \\
\textit{Question:} X'inhi l-età minima biex ikolli liċenzja sħiħa tal-karozza kategorija B? \\
\textit{Reference answer:} L-età minima għal liċenzja sħiħa kategorija B hi 18-il sena (learner minn 17). \\
\textit{Key facts (accepted aliases):} 18 (tmintax) \\
\textit{Expected citation:} S.L. 65.18
\par\vspace{5pt}
\noindent\textbf{MT12} \quad MT \quad Micromobility \quad hard \\
\textit{Question:} X'inhi l-veloċità massima ta' e-kickscooter fit-triq u fuq il-bankina? \\
\textit{Reference answer:} E-kickscooter jista' jagħmel sa 20 km/s fit-triq u 10 km/s fuq bankini u promenadi. \\
\textit{Key facts (accepted aliases):} 20 (għoxrin); 10 (għaxra); bankina (bankini, footpath, promenad, promenada, passaġġ, bankin) \\
\textit{Expected citation:} S.L. 65.32
\par\vspace{5pt}
\end{spacing}
